\subsection{RSA Algorithm}


\content{
        \begin{example} Compute the following: 
        \begin{enumerate}
        \item $ 3^9 \mod{15} $
        \item $ 2^9 \mod{15} $
        \item $ 9^9 \mod{15} $
        \item $14^9 \mod{15} $
        \end{enumerate}
        \end{example}
}{% presentation
\vspace{3in}
}{% nopresentation
\vspace{3in}
}{% comments
}

\content{
        \begin{example} What might we do with the above observation? 
        \end{example}
}{% presentation
\vspace{3in}
}{% nopresentation
\vspace{3in}
}{% comments
Factor 9 into 3 and 3, then notice that this gives rise to a 2-step process, we
could compute $b = a^3 \mod{15}$, then $c = b^3 \mod{15}$, but then
$c=a^9\mod{15} = a \mod{15}$. So we can recover what $a$ was. 
}

\content{
        \begin{definition} (RSA Algorithm) Each party wishing to communicate
        does the following:
        \begin{enumerate}
        \item Pick two prime numbers, $p$ and $q$. 
        \item Compute $n=p*q$.
        \item Compute $\phi(n) = (p-1)*(q-1)$
        \item Pick a public key $e$ satisfying $\gcd(e,\phi(n)) = 1$.
        \item Find $d$ such that $d*e \equiv 1\mod{\phi(n)}$.
        \item Publish $e$ and $n$ as public keys. 
        \item Guard $p$, $q$, and $\phi(n)$ as secret, and keeps $d$ secret as a
        secret key.
        \end{enumerate}
        \end{definition}

        \textbf{Question:} What makes this secure? What step is hard to compute
        without the secret keys? 
}{% presentation
\vspace{2in}
}{% nopresentation
\vspace{2in}
}{% comments
}

\content{
        \begin{example} Encrypt the word ``help'' using the RSA algorithm
        $n=3551$, $e=629$, and $d=1997$.
        \end{example}
}{% presentation
\vspace{4in}
}{% nopresentation
\vspace{4in}
}{% comments
}

\content{
        Ok... no more hand calculations for this lecture! Current best practices
        recommend a key length of at least 2048 bits, that means that $n$ should
        be a number with around 617 digits! Let's see a few examples with a
        computer.
}{% presentation
}{% nopresentation
}{% comments
Go to website, and see some examples of how this looks with larger primes. 
}
